% !PRINTED 277
\weberpage{277}
Eine Lösung dieser Kongruenz kann man immer aus einer Lösung der
Kongruenz (4) $\partial' c'' \equiv c \ (\mathrm{mod}\ a'')$ herleiten, indem
man zu $c''$ ein Vielfaches von $a''$ hinzufügt.

Die Kongruenz (7) hat dann nach § 69, (9) bis (13) zur Folge:
\weq{8}{\lambda \equiv n'\lambda'' \ (\mathrm{mod}\ 3),}
\weq{9}{
\begin{array}{ll}
\alpha \equiv \alpha''\partial'\partial_2', & \beta \equiv \beta''a'\partial_2' \\
\gamma \equiv \gamma''a'\partial_2', & \delta \equiv \delta''\partial'\partial_2'
\end{array}
\ (\mathrm{mod}\ 3).}

Es mögen nun $v', v_1', v_2'; v'', v_1'', v_2''$ dieselbe Bedeutung für die
Zahlen $n', n''$ haben, welche den $v, v_1, v_2$ in § 73, (3) für die Zahl
$n$ gegeben war, nämlich:
\[
v' = f\left(\frac{c' + \partial'\omega}{a'}\right), \qquad
v'' = f\left(\frac{c'' + \partial''\omega}{a''}\right),
\]
\weq{10}{
\begin{array}{ll}
v_1' = \left(\dfrac{2}{a'}\right) f_1\left(\dfrac{c' + \partial'\omega}{a'}\right), &
v_1'' = \left(\dfrac{2}{a''}\right) f_1\left(\dfrac{c'' + \partial''\omega}{a''}\right), \\[2ex]
v_2' = \left(\dfrac{2}{\partial'}\right) f_2\left(\dfrac{c' + \partial'\omega}{a'}\right), &
v_2'' = \left(\dfrac{2}{\partial''}\right) f_2\left(\dfrac{c'' + \partial''\omega}{a''}\right).
\end{array}
}

Wir wenden die Vertauschungstabelle (4), § 73 auf diese Funktionen an.
Da $n'$ unter allen Umständen durch 3 unteilbar ist, so sind die
kubischen Einheitswurzeln $\varrho', \sigma' = 1$ zu setzen, während
infolge der Kongruenzen (8), (9):
\weq{11}{\varrho'' = \varrho^{n'}, \quad \sigma'' = \sigma^{n'}}
wird. Wir erhalten hiernach folgende zusammengehörige Vertauschungen:
\weq{12}{
\begin{array}{llll}
\omega, & v', & v_1', & v_2', \\
-\dfrac{1}{\omega}, & v', & v_2', & v_1', \\
\omega+1, & e^{-\frac{n'\pi i}{24}}v_1', & e^{-\frac{n'\pi i}{24}}v', & e^{\frac{n'\pi i}{12}}v_2', \\
\hline
\omega, & v'', & v_1'', & v_2'', \\
-\dfrac{1}{\omega}, & \varrho^{n'}v'', & \varrho^{n'}v_2'', & \varrho^{n'}v_1'', \\
\omega+1, & \sigma^{n'}e^{-\frac{n''\pi i}{24}}v_1'', & \sigma^{n'}e^{-\frac{n''\pi i}{24}}v'', & \sigma^{n'}e^{\frac{n''\pi i}{12}}v_2'',
\end{array}
}
die zusammen mit den Vertauschungen (4), § 73 gelten.

Wir unterscheiden jetzt zwei Fälle.

1. Wenn
\weq{13}{(n'+1)(n''+1) = 8\mu \equiv 0 \ (\mathrm{mod}\ 8)}
%CONT>>
